\chapter{Number theory}

\section{Modular arithmetic}
	\kactlimport{ModularArithmetic.h}
	\kactlimport{ModInverseEuclid.h}
	\kactlimport{ModInverseSieve.h}
	\kactlimport{ModPow.h}
	\kactlimport{ModSum.h}
		\kactlimport{CntLattice.h}
	\kactlimport{ModMulLL.h}
	\kactlimport{ModLog.h}
	\kactlimport{ModSqrt.h}
	\kactlimport{PrimitiveRoot.h}
	\kactlimport{DiscreteRoot.h}


\section{Primality}
	\kactlimport{FastEratosthenes.h}
	\kactlimport{Eratosthenes.h}
	\kactlimport{SegmentedSieves.h}
	\kactlimport{Eratosthenes1e9.h}
	\kactlimport{PrimesToFactors.h}
	\kactlimport{MillerRabin.h}
	\kactlimport{PollardRho.h}

\section{Divisibility}
	\kactlimport{Euclid.h}
	\kactlimport{CRT.h}
	\kactlimport{LinearCong.h}
	\kactlimport{LinearDiophantine.h}

	\subsection{Number of solutions ax+by=c}
	General solution form is:
	$$
	x = x_0 + k \cdot \frac{b}{g}, \quad y = y_0 - k \cdot \frac{a}{g}
	$$
	We need $x > 0$ and $y > 0$.

	That gives constraints on $k$:
	$$k > \frac{-x_0 g}{b}, \quad k < \frac{y_0 g}{a}$$

	\subsection{Bézout's identity}
	For $a \neq $, $b \neq 0$, then $d=gcd(a,b)$ is the smallest positive integer for which there are integer solutions to
	$$ax+by=d$$
	If $(x,y)$ is one solution, then all solutions are given by
	$$\left(x+\frac{kb}{\gcd(a,b)}, y-\frac{ka}{\gcd(a,b)}\right), \quad k\in\mathbb{Z}$$

	\kactlimport{phiFunction.h}

\section{Fractions}
	\kactlimport{ContinuedFractions.h}
	\kactlimport{FracBinarySearch.h}

\section{Pythagorean Triples}
	The Pythagorean triples are uniquely generated by
	\[ a=k\cdot (m^{2}-n^{2}),\ \,b=k\cdot (2mn),\ \,c=k\cdot (m^{2}+n^{2}), \]
	with $m > n > 0$, $k > 0$, $m \bot n$ (coprime), and either $m$ or $n$ even.

\section{Primes}
	$p=962592769$ is such that $2^{21} \mid p-1$, which may be useful. For hashing
	use 970592641 (31-bit number), 31443539979727 (45-bit), 3006703054056749
	(52-bit). There are 78498 primes less than 1\,000\,000.

	Primitive roots exist modulo any prime power $p^a$, except for $p = 2, a > 2$, and there are $\phi(\phi(p^a))$ many.
	For $p = 2, a > 2$, the group $\mathbb Z_{2^a}^\times$ is instead isomorphic to $\mathbb Z_2 \times \mathbb Z_{2^{a-2}}$.

\section{Estimates}
	\textbf{Divisor Bound}: The number of divisors $d(n)$ grows slowly.
	\[ \sum_{d|n} d = O(n \log \log n) \]

	The number of divisors of $n$ is at most around 100 for $n < 5e4$, 500 for $n < 1e7$, 1344 for $n < 1e9$, 2000 for $n < 1e10$, 103680 for $n < 1e18$, 200\,000 for $n < 1e19$.
	Approx bound: $d(n) \approx n^{1.066/\ln \ln n}$.

\section{Mobius Function}
	\kactlimport{Mobius.h}
\[
	\mu(n) = \begin{cases} 0 & n \textrm{ is not square free}\\ 1 & n \textrm{ has even number of prime factors}\\ -1 & n \textrm{ has odd number of prime factors}\\\end{cases}
\]
  Mobius Inversion:
  \[ g(n) = \sum_{d|n} f(d) \Leftrightarrow f(n) = \sum_{d|n} \mu(d)g(n/d) \]

  Other useful formulas/forms:

  \[ \sum_{d | n} \mu(d) = [ n = 1] \]

  \[ g(n) = \sum_{n|d} f(d) \Leftrightarrow f(n) = \sum_{n|d} \mu(d/n)g(d) \]

 \[ g(n) = \sum_{1 \leq m \leq n} f(\left\lfloor\frac{n}{m}\right \rfloor ) \Leftrightarrow f(n) = \sum_{1\leq m\leq n} \mu(m)g(\left\lfloor\frac{n}{m}\right\rfloor) \]

 \section{Mega Sieve}
	\kactlimport{MegaSieve.h}

\section{Theorems}
If $\gcd(A, B) = 1$, then $\gcd(A+B, A \cdot B) = 1$
\subsection{Modular \& Factorial}
	\kactlimport{Legendre.h}
	\textbf{Wilson's Theorem}: \\ $(p-1)! \equiv -1 \pmod p$ iff $p$ is prime. Any integer $n > 1$ is prime iff $(n-1)! + 1$ is divisible by $n$. \\
	\textbf{Legendre's Formula}: \\ The exponent of prime $p$ in $n!$ is $\sum_{k=1}^{\infty} \lfloor \frac{n}{p^k} \rfloor = \frac{n - S_p(n)}{p-1}$, where $S_p(n)$ is the sum of digits of $n$ in base $p$. \\
	\textbf{Lifting The Exponent (LTE)}: \\ If $p \mid x-y$, $p \nmid n, x, y$, then $v_p(x^n-y^n) = v_p(x-y) + v_p(n)$. If $p=2$, $v_2(x^n-y^n) = v_2(x-y) + v_2(n) + v_2(x+y) - 1$ for even $n$. \\
\subsection{Sums of Squares}
	\textbf{Fermat (2 Squares)}: \\ $n = a^2+b^2$ iff $v_p(n)$ is even for every prime $p \equiv 3 \pmod 4$. \\
	\textbf{Legendre (3 Squares)}: \\ $n = a^2+b^2+c^2$ iff $n \neq 4^k(8m+7)$. \\
	\textbf{Lagrange (4 Squares)}: \\ Every $n \in \mathbb{N}$ is the sum of 4 integer squares. \\
\subsection{Diophantine \& Sequences}
	\textbf{Chicken McNugget}: \\ For coprime $a, b$, the largest integer that cannot be written as $ax+by$ ($x,y \ge 0$) is $ab - a - b$. There are exactly $\frac{(a-1)(b-1)}{2}$ such integers. \\
	\textbf{Zeckendorf}: \\ Every positive integer can be uniquely represented as a sum of non-consecutive Fibonacci numbers ($F_1=1, F_2=2, \dots$). \\
	\textbf{Bertrand's Postulate}: \\ For $n > 1$, there is always a prime $p$ such that $n < p < 2n$. \\
	\textbf{Dirichlet}: \\ If $\gcd(a, b)=1$, the sequence $an+b$ contains infinitely many primes. \\
\subsection{Conjectures}
	\textbf{Goldbach (Strong)}: \\ Every even integer $n > 2$ is the sum of two primes. \\
	\textbf{Goldbach (Weak)}: \\ Every odd integer $n > 5$ is the sum of three primes. \\